\chapter{Robustifying \texttt{DMPlex} for Naturally Mixed Element Families}
\label{ch:dmplex-mixed-topology}

\section{Motivation}

Up to this point, the \texttt{Domain}--\texttt{Model} pipeline relied on a
very restrictive interpretation of \texttt{DMPlex}: an uninterpolated
cell--vertex graph where every analysis element was treated as a ``cell'' and
the code assumed that querying height~$1$ produced the vertices. That
assumption is brittle even for standard meshes and becomes conceptually wrong
as soon as one wants to mix, in the same discrete domain, elements of
different topological dimension such as:
\begin{itemize}
  \item beams ($1$D topology embedded in $\mathbb{R}^3$),
  \item shells ($2$D topology embedded in $\mathbb{R}^3$),
  \item continua ($3$D topology embedded in $\mathbb{R}^3$).
\end{itemize}

The immediate objective of this refactor is therefore modest but essential:
\emph{keep the current non-interpolated DAG for performance and simplicity,
while making the \texttt{DMPlex} object semantically correct enough to host
mixed element families naturally.}

\section{Design Decision}

The adopted strategy is intentionally incremental:
\begin{enumerate}
  \item Keep the DAG non-interpolated. Cells still connect directly to
        vertices; explicit edges and faces are not materialised yet.
  \item Encode the missing topological semantics using \texttt{DMLabel}s.
  \item Drive the algebraic layout from the only invariant that really matters
        for the current solver: analysis unknowns live on depth-$0$ vertices.
\end{enumerate}

This preserves the low-cost setup path and avoids polluting the hot assembly
loops with new runtime indirections. The extra work happens only once during
mesh/DMPlex setup.

\section{Implemented Semantics}

Three labels are now attached to the \texttt{DMPlex} points generated by
\texttt{Domain::assemble\_sieve()}:
\begin{description}
  \item[\texttt{fall\_n.point\_role}] distinguishes between vertex points and
        cell points. The current values are:
        \begin{itemize}
          \item $0$ for vertices,
          \item $1$ for cells.
        \end{itemize}
  \item[\texttt{fall\_n.topological\_dim}] stores the intrinsic topological
        dimension of each point. Vertices receive value $0$, while cells
        receive the topological dimension reported by
        \texttt{ElementGeometry::topological\_dimension()}.
  \item[\texttt{fall\_n.physical\_group\_id}] stores an integer identifier for
        the physical group carried by a geometric element, when such group
        exists.
\end{description}

These labels are cheap, explicit, and stable. More importantly, they separate
two different concerns that had previously been conflated:
\begin{itemize}
  \item the \emph{graph} structure of the mesh;
  \item the \emph{topological meaning} of each graph point.
\end{itemize}

\section{Why \texttt{Depth} and Not \texttt{Height}}

The previous implementation used:
\begin{verbatim}
DMPlexGetHeightStratum(dm, 1, &vStart, &vEnd);
\end{verbatim}
and interpreted that stratum as the vertices. This is only safe for a very
specific cell--vertex graph and becomes false once the mesh is interpolated or
its topology is enriched.

The robust statement is:
\begin{quote}
analysis vertices are the depth-$0$ points of the \texttt{DMPlex}.
\end{quote}

Accordingly, \texttt{Model::set\_sieve\_layout()} was rewritten to define the
section over the full chart while conceptually placing DOFs only on depth-$0$
points. Cell points keep zero DOFs.

\section{Dimension of the \texttt{DMPlex}}

Another correction concerns the dimension assigned to the \texttt{DM}. The old
code forced:
\begin{verbatim}
DMSetDimension(dm, dim);
\end{verbatim}
where \texttt{dim} is the embedding space dimension of the model, not the
intrinsic dimension of the discrete topology.

This is now replaced by:
\begin{equation}
\dim(\texttt{DM}) = \max_{e \in \mathcal{T}_h} \operatorname{topodim}(e),
\end{equation}
that is, the maximum topological dimension among the stored elements. For a
mixed beam--shell model embedded in $\mathbb{R}^3$, the resulting
\texttt{DMPlex} has dimension~$2$; for a mixed beam--shell--solid model, it has
dimension~$3$.

\section{Compile-Time and HPC Considerations}

This refactor was designed to preserve the original performance intent:
\begin{itemize}
  \item Geometry, interpolation, and quadrature remain statically typed.
  \item No new virtual dispatch was introduced inside numerical integration,
        constitutive update, or element assembly loops.
  \item The new metadata is written once during \texttt{assemble\_sieve()} and
        then consumed through PETSc's own label machinery.
  \item The DAG itself is still cell--vertex, which keeps setup lighter than a
        fully interpolated mesh and reduces memory pressure.
\end{itemize}

In other words, the refactor improves semantic correctness of the mesh layer
without paying for it in the hot path.

\section{Validation}

A dedicated regression test was added:
\begin{center}
  \texttt{tests/test\_dmplex\_mixed\_topology.cpp}
\end{center}
It builds a small domain containing:
\begin{itemize}
  \item one $1$D beam geometry,
  \item one $2$D shell geometry,
  \item shared depth-$0$ vertices.
\end{itemize}

The test verifies:
\begin{enumerate}
  \item the \texttt{DM} dimension is the maximum topological dimension;
  \item height-$0$ contains the two cells and depth-$0$ the four vertices;
  \item the three labels above are present and correctly populated;
  \item a polymorphic structural \texttt{Model} assigns DOFs only to vertices,
        leaving cell points with zero DOFs.
\end{enumerate}

\section{Current Scope and Next Step}

This chapter documents a necessary intermediate stage, not the final one.
Mixed families now coexist naturally in the same \texttt{DMPlex} at the
semantic level, but the mesh is still not fully interpolated. Therefore:
\begin{itemize}
  \item the solver can reason correctly about vertices versus cells;
  \item element families of different intrinsic dimension can share the same
        \texttt{Domain} and \texttt{DMPlex};
  \item future work may still introduce explicit edges/faces when the library
        needs richer topological traversal or native mixed-dimensional coupling
        operators.
\end{itemize}

This staging is intentional. It respects the architectural principle adopted
throughout \texttt{fall\_n}: only introduce abstractions that correspond to a
real numerical need, and keep the compile-time structure of the kernels as
static as possible.
